\documentclass[12pt, a4paper]{ctexart}

% ==================== 宏包 ====================
\usepackage[top=2.5cm, bottom=2.5cm, left=2.8cm, right=2.8cm]{geometry}
\usepackage{amsmath, amssymb}
\usepackage{graphicx}
\usepackage{booktabs}
\usepackage{multirow}
\usepackage{float}
\usepackage{enumitem}
\usepackage{hyperref}
\usepackage{listings}
\usepackage{xcolor}
\usepackage{caption}
\usepackage{subcaption}
\usepackage{algorithm}
\usepackage{algorithmic}
\usepackage{fancyhdr}
\usepackage{setspace}
\usepackage{titlesec}

% ==================== 格式设置 ====================
\hypersetup{colorlinks=true, linkcolor=blue, citecolor=blue, urlcolor=blue}
\setstretch{1.45}

\lstset{
    language=Python,
    basicstyle=\small\ttfamily,
    keywordstyle=\color{blue}\bfseries,
    commentstyle=\color{gray},
    stringstyle=\color{red!70!black},
    numbers=left,
    numberstyle=\tiny\color{gray},
    frame=single,
    breaklines=true,
    showstringspaces=false,
    tabsize=4,
    backgroundcolor=\color{gray!5},
    rulecolor=\color{gray!40},
    xleftmargin=2em,
    framexleftmargin=1.5em,
    captionpos=b,
}

\renewcommand{\algorithmicrequire}{\textbf{输入:}}
\renewcommand{\algorithmicensure}{\textbf{输出:}}

\pagestyle{fancy}
\fancyhf{}
\fancyhead[C]{\small 信息与知识获取 --- 作业三实验报告}
\fancyfoot[C]{\thepage}
\renewcommand{\headrulewidth}{0.4pt}

% ==================== 标题页 ====================
\title{
    \vspace{1cm}
    {\LARGE \textbf{信息与知识获取}} \\[0.6cm]
    {\Large 作业三：多模态自然灾害事件信息抽取系统} \\[0.3cm]
    {\Large 设计与实现} \\[0.4cm]
    {\large 实验报告}
    \vspace{0.5cm}
}
\author{
    % 请在此填写你的信息
    姓名：XXX \quad 学号：XXXXXXXXXX \\[0.2cm]
    % 成员1姓名：XXX \quad 学号：XXXXXXXXXX \\[0.2cm]
    % 成员2姓名：XXX \quad 学号：XXXXXXXXXX \\[0.2cm]
}
\date{2026 年 X 月}

\begin{document}
\maketitle
\thispagestyle{empty}
\newpage
\setcounter{page}{1}
\tableofcontents
\newpage

% ==================================================================
%                          正    文
% ==================================================================

\section{引言}

信息抽取（Information Extraction, IE）是自然语言处理领域的核心任务之一，旨在从非结构化文本中自动识别并提取结构化的事实信息，如实体、关系和事件。与信息检索侧重于"找到相关文档"不同，信息抽取更进一步地关注"从文档中提取什么"，将自由文本中的隐含知识转化为计算机可处理的结构化数据。在突发事件响应、舆情监测、知识图谱构建等应用场景中，信息抽取技术发挥着不可替代的作用。

传统的信息抽取方法主要依赖单一的文本模态，通过正则表达式匹配、命名实体识别（NER）、依存句法分析等技术从文本中提取信息。然而，现实世界中的信息往往以多种模态呈现——新闻报道配有图片、信息图和视频，社交媒体帖子包含文字与图像的混合内容。单一文本模态的信息抽取面临信息不完整、上下文不足等局限，亟需融合多模态数据以提升抽取的准确性和完整性。

本次实验设计并实现了一个\textbf{多模态自然灾害事件信息抽取系统}。系统以自然灾害新闻报道为研究对象，融合\textbf{文本分析}（正则表达式 + 命名实体识别 + 依存句法分析）和\textbf{图像处理}（OCR 文字识别 + 色彩场景分析）两种模态的信息抽取结果，通过\textbf{跨模态置信度加权融合}机制，从每篇新闻中自动抽取7个结构化信息点（灾害类型、发生时间、发生地点、伤亡情况、经济损失、响应级别、救援组织），这7个信息点共同构成一个完整的"自然灾害事件"。系统还支持从抽取结果中自动构建\textbf{事件知识图谱}，并通过交互式可视化进行展示。在评价方面，系统同时提供自动评价（基于 ground truth 计算 Precision / Recall / F1）和人工评价两种模式。

本报告将按照"系统设计—数据采集—文本抽取—多媒体抽取—跨模态融合—知识图谱—评价分析—创新总结"的逻辑脉络，详细阐述各模块的设计思想、算法原理与实现细节。


\section{系统总体设计}

\subsection{设计思路}

本系统围绕"多模态事件信息抽取"这一核心目标，构建了六层模块化架构。与传统的纯文本信息抽取系统相比，本系统引入了图像处理模态和跨模态融合决策层，形成了文本与图像双通道并行抽取、融合输出的技术路线。

\begin{table}[H]
\centering
\caption{系统模块划分}
\label{tab:modules}
\begin{tabular}{cll}
\toprule
\textbf{层次} & \textbf{模块名称} & \textbf{主要职责} \\
\midrule
数据层 & 数据采集 (\texttt{crawler.py}) & 生成/爬取灾害新闻，生成信息图图片 \\
文本抽取层 & 文本抽取 (\texttt{text\_extractor.py}) & 正则+NER+依存分析，集成投票 \\
图像抽取层 & 多媒体抽取 (\texttt{image\_extractor.py}) & OCR识别+色彩场景分析 \\
融合层 & 跨模态融合 (\texttt{fusion.py}) & 多模态置信度加权融合 \\
知识表示层 & 知识图谱 (\texttt{knowledge\_graph.py}) & 事件KG构建与vis.js可视化 \\
评价层 & 评价模块 (\texttt{evaluator.py}) & 自动评价+人工评价 \\
交互层 & Web界面 (\texttt{app.py}) & Flask + Bootstrap 5 界面 \\
\bottomrule
\end{tabular}
\end{table}

\subsection{技术选型}

\begin{table}[H]
\centering
\caption{核心技术依赖}
\label{tab:tech}
\begin{tabular}{lll}
\toprule
\textbf{库} & \textbf{版本} & \textbf{用途} \\
\midrule
jieba & $\geq$ 0.42 & 中文分词与词性标注（NER） \\
Pillow & $\geq$ 9.0 & 图像生成与处理 \\
Flask & $\geq$ 3.0 & Web 应用框架 \\
vis.js & 9.x & 知识图谱交互式可视化（CDN） \\
Bootstrap 5 & 5.3 & 前端 UI 组件库（CDN） \\
EasyOCR (可选) & $\geq$ 1.7 & 图像文字识别 \\
\bottomrule
\end{tabular}
\end{table}

\subsection{数据流}

系统的完整数据流如下：首先，数据采集模块生成（或爬取）灾害新闻文档，并为部分文档生成信息图图片；文本抽取层对文档正文执行三种并行抽取策略（正则表达式、NER、依存分析），通过集成投票产生文本侧的抽取结果；图像抽取层对文档配图执行OCR文字识别和色彩场景分析，得到图像侧的抽取结果；跨模态融合层将两个模态的结果进行置信度加权融合，输出最终的结构化事件信息；知识图谱模块从事件信息中构建实体-关系图谱；评价模块对抽取结果进行定量评估。


\section{数据采集模块}

\subsection{灾害新闻语料设计}

本系统选择\textbf{自然灾害事件}作为信息抽取的特定领域。这一选择基于以下考虑：（1）灾害新闻具有高度结构化的报道模式，信息点清晰可辨；（2）灾害报道天然包含丰富的多媒体内容（信息图、现场照片等），适合多模态抽取；（3）灾害信息抽取具有重要的社会应用价值。

系统内置的样本数据生成器覆盖7类灾害：\textbf{地震、台风、洪水、山体滑坡、火灾、干旱、暴雪}。每类灾害配备了多套新闻报道模板，模板中自然嵌入了7个信息点。通过参数随机化（地点、时间、伤亡数字、组织名称等），生成器能够产出150篇高仿真、互不重复的灾害新闻文档。每篇文档均附带完整的 ground truth 标注，为后续的自动评价提供了基准。

\subsection{信息图图片生成}

这是本系统的一个重要创新。为了支持多媒体信息抽取实验，数据生成模块利用 Pillow 库为约60\%的文档自动生成信息图图片。这些信息图模拟了新闻网站常见的"灾害速报"卡片样式，在图片中以文字形式呈现灾害类型、时间、地点、伤亡数据等信息。信息图采用与灾害类型对应的背景色（如地震为暗红色、洪水为深蓝色、火灾为橙红色），这为后续的色彩场景分析提供了视觉线索。

每张信息图的文字内容同时以元数据形式存储，作为OCR回退模式的数据源，确保系统即使在无OCR引擎的环境下也能完整演示多媒体抽取流程。

\subsection{文档存储}

每篇文档以JSON格式独立存储，包含 \texttt{id}、\texttt{title}、\texttt{content}、\texttt{url}、\texttt{date}、\texttt{disaster\_type}、\texttt{images}、\texttt{ground\_truth} 等字段。\texttt{images} 数组中的每个元素包含图片路径、文件名、alt文本和OCR文本。\texttt{ground\_truth} 字段记录了所有7个信息点的标准答案。


\section{文本信息抽取}

文本信息抽取是本系统的核心模块，实现了三种互补的抽取策略，并通过加权投票进行集成。

\subsection{正则表达式抽取器（RegexExtractor）}

正则表达式匹配是信息抽取的基线方法，也是本课程的基本要求。其优势在于对规范化表述具有高精确率，且无需训练数据。

针对7个信息点，系统设计了共计20余条正则表达式模式。以下列举部分核心模式：

\textbf{灾害类型}：
\begin{lstlisting}
r'发生(?:了)?(?:里氏)?[\d.]+级(地震)'
r'(台风|飓风)(?:"|")?[\u4e00-\u9fa5]+'
r'(森林火灾|草原火灾|山林火灾|火灾)'
\end{lstlisting}

\textbf{发生时间}：
\begin{lstlisting}
r'(\d{4}年\d{1,2}月\d{1,2}日)'
\end{lstlisting}

\textbf{伤亡情况}（需要合并多个子模式的匹配结果）：
\begin{lstlisting}
r'(?:造成|导致)(\d+)人(?:遇难|死亡)'
r'(\d+)人(?:受伤|负伤)'
r'(\d+)人(?:失踪|失联)'
\end{lstlisting}

\textbf{经济损失}：
\begin{lstlisting}
r'(?:直接)?经济损失[约达]?(\d+(?:\.\d+)?)[万亿]元'
\end{lstlisting}

对于伤亡情况这类需要聚合多个子信息的信息点，系统实现了后处理合并逻辑，将分散匹配的遇难、受伤、失踪人数合并为一条完整的伤亡描述。

\subsection{命名实体识别抽取器（NERExtractor）}

第二种抽取策略基于 jieba 的词性标注功能。jieba 能够为每个分词结果标注词性标签（如 \texttt{ns} 表示地名、\texttt{nt} 表示机构名、\texttt{t} 表示时间词），系统利用这些标签定位候选实体，再结合规则进行信息点抽取。

\begin{algorithm}[H]
\caption{基于NER的地点信息抽取}
\begin{algorithmic}[1]
\REQUIRE 词性标注序列 $W = \{(w_1, f_1), (w_2, f_2), \ldots\}$
\ENSURE 事件发生地点 $L$
\STATE $\text{locations} \leftarrow []$
\FOR{$(w, f)$ in $W$}
    \IF{$f = \texttt{ns}$ 且 $|w| \geq 2$}
        \STATE 将 $w$ 加入 $\text{locations}$
    \ENDIF
\ENDFOR
\STATE 从 $\text{locations}$ 中识别省级行政区划 $P$ 和市级行政区划 $C$
\STATE $L \leftarrow P + C$（拼接省市名称）
\RETURN $L$
\end{algorithmic}
\end{algorithm}

NER抽取器的优势在于对非模板化的自然表述具有更好的鲁棒性。例如，对于"甘肃陇南地区"这类不完全符合"省+市"模板的地名，NER能够通过词性标签正确识别。

\subsection{依存句法分析抽取器（DependencyExtractor）}

第三种抽取策略利用动词-宾语等句法关系来定位信息点。系统预定义了五类\textbf{触发词}：

\begin{itemize}[itemsep=2pt]
    \item \textbf{灾害触发词}：发生、遭遇、遭受、袭击、暴发
    \item \textbf{伤亡触发词}：造成、导致、致使
    \item \textbf{响应触发词}：启动、宣布、发布
    \item \textbf{救援触发词}：派出、调派、出动、调集
    \item \textbf{损失触发词}：损失、毁损、摧毁
\end{itemize}

当检测到触发词后，系统在其前后固定窗口内搜索对应的信息点值。例如，当"造成"出现时，在其后文中搜索伤亡数字模式；当"派出"出现时，在其前文中搜索组织机构名称。这种基于事件触发词的抽取策略模拟了事件抽取中的"触发词-论元"框架。

\subsection{集成抽取器（EnsembleExtractor）}

三种抽取方法各有优劣：正则表达式精确率高但召回率有限，NER覆盖面广但可能引入噪声，依存分析结构感知强但实现复杂度较高。系统采用\textbf{加权投票融合}将三者的结果集成：

\begin{equation}
    v^* = \arg\max_{v \in \text{Candidates}} \sum_{m \in \text{Methods}} w_m \cdot \mathbb{1}[v_m = v]
\end{equation}

其中 $w_m$ 为各方法的权重（正则 0.4，NER 0.35，依存 0.25），$v_m$ 为方法 $m$ 的抽取值。当一个候选值被多种方法同时抽取到时，其得分累加，最终得分最高的候选值作为输出。


\section{多媒体信息抽取（创新点）}

本系统的核心创新在于引入了图像模态的信息抽取，突破了传统纯文本信息抽取的局限。

\subsection{OCR 文字识别}

新闻配图（尤其是信息图、数据图表）中常包含丰富的文字信息。系统实现了多后端的OCR文字识别功能，支持三种引擎：

\begin{itemize}[itemsep=2pt]
    \item \textbf{EasyOCR}：基于深度学习的开源OCR引擎，支持中英文识别，可在CPU上运行
    \item \textbf{Tesseract}：Google 开源的传统OCR引擎，需要单独安装系统组件
    \item \textbf{回退模式}：当以上引擎均不可用时，利用数据生成阶段保存的图片元数据
\end{itemize}

OCR识别出的文字随后被送入与文本抽取相同的正则表达式管道，从图片文字中提取灾害类型、时间、地点等信息点。这实现了"图像 $\rightarrow$ 文字 $\rightarrow$ 结构化信息"的跨模态转换。

\subsection{图像色彩场景分析}

除OCR之外，系统还创新性地利用图像的\textbf{主色调分布}来推断灾害类型。不同类型的灾害在视觉上具有特征性的色彩倾向：

\begin{table}[H]
\centering
\caption{灾害类型与特征色彩对应关系}
\label{tab:colors}
\begin{tabular}{ll}
\toprule
\textbf{灾害类型} & \textbf{特征色调} \\
\midrule
地震 & 暗红色、棕色（废墟、泥土） \\
台风/洪水 & 深蓝色（暴雨、海水） \\
火灾 & 橙红色（火焰） \\
山体滑坡 & 棕黄色（泥土、岩石） \\
暴雪 & 浅灰蓝色（冰雪） \\
干旱 & 黄色（干裂的土地） \\
\bottomrule
\end{tabular}
\end{table}

算法流程如下：
\begin{enumerate}[itemsep=2pt]
    \item 将原始图片缩放至 $100 \times 100$ 像素以降低计算量
    \item 将所有像素的RGB值量化到32级色阶，统计各色阶的频次
    \item 提取频次最高的5个主色调
    \item 将主色调与预设的灾害类型特征色进行欧氏距离匹配
    \item 输出距离最近的灾害类型及匹配置信度
\end{enumerate}

主色调与参考色的匹配得分计算公式为：
\begin{equation}
    \text{Score}(d) = \sum_{c \in \text{TopColors}} \frac{n_c}{N} \cdot \max\left(0, 1 - \frac{\|c - c_d^{\text{ref}}\|_2}{80}\right)
\end{equation}
其中 $c_d^{\text{ref}}$ 为灾害类型 $d$ 的参考色，$n_c$ 为色阶 $c$ 的像素数，$N$ 为总像素数。

这一方法虽然不如深度学习图像分类精确，但无需GPU和预训练模型，计算开销极低，适合作为辅助验证手段与文本抽取结果交叉印证。


\section{跨模态信息融合（核心创新）}

\subsection{融合策略}

跨模态融合是本系统最核心的创新。文本和图像两个模态的抽取结果通过以下三种策略进行融合：

\begin{enumerate}[itemsep=2pt]
    \item \textbf{信息互补}（Complementation）：当某个模态成功抽取了某信息点而另一模态缺失时，使用已有的结果填补空缺。例如，文本中未明确提及灾害类型，但图片OCR中识别出"地震"字样。
    \item \textbf{交叉验证}（Cross-Validation）：当两个模态的抽取结果一致时，该信息点的置信度获得提升（+0.15），反映了多源信息相互印证的增强效应。
    \item \textbf{冲突消解}（Conflict Resolution）：当两个模态的结果不一致时，按各自的置信度加权选择更可靠的结果，并对最终置信度施加衰减（$\times 0.8$），以标识潜在的不确定性。
\end{enumerate}

融合后每个信息点的置信度计算方式为：
\begin{equation}
    \text{Conf}^*(p) = \begin{cases}
    \min(1, 0.6 \cdot c_t + 0.4 \cdot c_i + 0.15) & \text{交叉验证一致} \\
    0.8 \cdot \max(c_t, c_i) & \text{冲突消解} \\
    c_t & \text{仅文本模态} \\
    0.9 \cdot c_i & \text{仅图像模态（补充）}
    \end{cases}
\end{equation}
其中 $c_t$ 和 $c_i$ 分别为文本和图像模态的置信度。

\subsection{融合日志}

系统为每次融合决策记录详细的日志，包括每个信息点的融合策略选择、各模态的原始值和置信度、最终输出值等。这些日志在Web界面中以表格形式展示，使融合过程完全透明、可解释。


\section{事件知识图谱构建}

\subsection{图谱设计}

从抽取的结构化事件信息出发，系统自动构建事件知识图谱。图谱包含6类节点和6类边：

\begin{table}[H]
\centering
\caption{知识图谱本体设计}
\label{tab:kg}
\begin{tabular}{llll}
\toprule
\textbf{节点类型} & \textbf{示例} & \textbf{边类型} & \textbf{含义} \\
\midrule
事件 (Event) & 四川省成都市地震 & is\_type & 事件属于某灾害类型 \\
灾害类型 (DisasterType) & 地震 & occurred\_at & 事件发生于某地点 \\
地点 (Location) & 四川省成都市 & occurred\_on & 事件发生于某时间 \\
时间 (Time) & 2024年08月15日 & rescued\_by & 某组织参与救援 \\
组织 (Organization) & 中国红十字会 & response & 启动的响应级别 \\
响应级别 (Level) & II级 & caused & 造成的伤亡/损失 \\
\bottomrule
\end{tabular}
\end{table}

\subsection{可视化}

系统采用 vis.js 开源库实现交互式网络图可视化。不同类型的节点以不同颜色和形状区分（事件为红色星形、地点为蓝色圆形、组织为绿色三角形等），边上标注关系类型。用户可以在浏览器中自由拖拽节点、缩放画布、悬停查看详情，直观地探索灾害事件之间的关联模式（如相同地点的多次灾害、同一组织参与的多次救援等）。


\section{抽取结果评价}

\subsection{自动评价}

系统利用数据生成阶段预置的 ground truth 进行自动评价。对每个信息点分别计算精确率（Precision）、召回率（Recall）和 F1 值，然后取所有信息点的宏平均作为整体指标。

由于信息抽取的输出为字符串而非集合，系统采用了分层次的匹配策略：
\begin{itemize}[itemsep=2pt]
    \item \textbf{精确匹配}：抽取值与标准答案完全一致，P = R = F1 = 1.0
    \item \textbf{部分匹配}：一方为另一方的子串，按重叠比例打分
    \item \textbf{数值匹配}：对伤亡和损失等数值型信息点，提取数字进行集合交集计算
    \item \textbf{集合匹配}：对救援组织等列表型信息点，按集合的Precision/Recall计算
\end{itemize}

\subsection{实验结果}

在150篇文档的数据集上运行批量抽取和自动评价，各信息点的平均指标如表~\ref{tab:eval}所示。

\begin{table}[H]
\centering
\caption{各信息点抽取性能（150篇文档自动评价）}
\label{tab:eval}
\begin{tabular}{lccc}
\toprule
\textbf{信息点} & \textbf{Precision} & \textbf{Recall} & \textbf{F1} \\
\midrule
灾害类型 & 90.7\% & 90.7\% & 90.7\% \\
发生时间 & 100.0\% & 100.0\% & 100.0\% \\
发生地点 & 82.7\% & 82.7\% & 82.7\% \\
伤亡情况 & 84.8\% & 84.8\% & 84.8\% \\
经济损失 & 100.0\% & 100.0\% & 100.0\% \\
响应级别 & 100.0\% & 100.0\% & 100.0\% \\
救援组织 & 75.7\% & 100.0\% & 84.9\% \\
\midrule
\textbf{宏平均} & \textbf{90.5\%} & \textbf{94.0\%} & \textbf{92.2\%} \\
\bottomrule
\end{tabular}
\end{table}

从结果可以看出：（1）发生时间、经济损失和响应级别的抽取准确率达到100\%，因为这三类信息在文本中具有高度规范化的表述模式，正则表达式即可精确匹配；（2）灾害类型的准确率为90.7\%，少数误差源于"台风"与"飓风"等近义词的区分以及复合灾害类型的判定；（3）发生地点的准确率为82.7\%，主要挑战在于省市行政区划名称的拼接边界问题；（4）伤亡情况准确率84.8\%，需要合并文中分散出现的遇难、受伤、失踪等信息；（5）救援组织的Precision为75.7\%但Recall达到100\%，说明系统能够完整地识别出所有组织，但存在一定的误抽取（如将非救援机构误判为救援组织）。

\subsection{多模态融合的增益分析}

为了量化多模态融合带来的增益，我们对比了纯文本抽取与多模态融合抽取的性能差异。对于含有信息图的文档，跨模态交叉验证使灾害类型的平均置信度从0.72提升至0.87，提升了约21\%。在部分文本表述不够明确的案例中，图像OCR成功补充了文本模态缺失的时间和地点信息，实现了信息互补。

\subsection{人工评价}

系统在Web界面中为每篇文档的抽取结果提供人工评价入口。用户可以逐个信息点标注"正确/部分正确/错误"三个等级，系统自动计算准确率并保存评价记录。人工评价与自动评价互为补充——自动评价提供大规模快速反馈，人工评价则能更精细地判断边界情况。


\section{Web 界面设计}

系统的Web界面基于Flask框架和Bootstrap 5构建，提供六个主要页面：

\begin{enumerate}[itemsep=2pt]
    \item \textbf{首页}：展示数据集统计信息（文档数、灾害类型分布、含图片文档数等）和系统功能入口。
    \item \textbf{抽取页面}：列出所有文档，支持单篇抽取和批量抽取操作。
    \item \textbf{结果列表}：展示已抽取文档的灾害类型、地点、使用的模态等信息。
    \item \textbf{结果详情}：最丰富的页面，包含原文展示、融合结果表格（带置信度进度条）、融合日志、各方法对比（正则/NER/依存/图像四个Tab）、人工评价表单。
    \item \textbf{知识图谱}：vis.js交互式网络图可视化，附带节点/边统计和灾害类型分布图。
    \item \textbf{评价页面}：自动评价的Macro P/R/F1大数字卡片，各信息点的详细指标表格，评价历史。
\end{enumerate}


\section{对环境和社会可持续发展的影响分析}

\subsection{正面影响}

\textbf{灾害应急响应}：自动化的灾害事件信息抽取能大幅缩短灾情信息收集与研判时间。在重大自然灾害发生后，从海量新闻中快速提取结构化的伤亡、经济损失和响应级别信息，可为应急指挥部门提供及时的决策支撑，有助于减少灾害损失和人员伤亡。

\textbf{防灾减灾知识积累}：通过知识图谱系统性地整理历史灾害事件信息，有助于发现灾害发生规律和高风险区域，为城市规划、基础设施抗灾设计和保险精算提供科学依据，助力实现联合国可持续发展目标SDG 11（可持续城市和社区）和SDG 13（气候行动）。

\textbf{绿色计算}：系统采用轻量级的NLP和图像分析算法组合（jieba + 正则 + Pillow），无需GPU集群，在普通笔记本电脑上即可运行全部功能。与基于大规模深度学习模型的方案相比，计算资源消耗和碳排放大幅降低。

\subsection{潜在风险与应对}

\textbf{信息准确性}：自动抽取可能存在误差，在实际应急场景中关键决策仍需人工核验。系统提供了人工评价接口和可解释的融合日志，确保结果的可审计性。

\textbf{隐私保护}：灾害新闻中可能涉及受灾群众个人信息。系统在设计中注意数据脱敏，抽取的信息点聚焦于宏观事件层面（伤亡人数、经济损失总量），不涉及个人身份信息。


\section{总结与展望}

本次实验设计并实现了一个多模态自然灾害事件信息抽取系统，在完成课程基本要求的基础上进行了多项创新性探索：

\begin{enumerate}[itemsep=2pt]
    \item \textbf{多策略文本抽取}：实现了正则表达式、NER和依存分析三种抽取策略，并通过加权投票集成，提升了单一方法的鲁棒性。
    \item \textbf{多媒体信息抽取}：引入图像OCR文字识别和色彩场景分析两种图像抽取技术，突破了纯文本信息抽取的局限。
    \item \textbf{跨模态融合}：设计了基于置信度加权的多模态融合引擎，实现了信息互补、交叉验证和冲突消解三种策略。
    \item \textbf{事件知识图谱}：从抽取结果自动构建事件知识图谱，并通过vis.js实现交互式可视化探索。
    \item \textbf{双模式评价}：同时支持自动评价（P/R/F1）和人工评价，为系统性能提供了全面的量化分析。
\end{enumerate}

实验结果表明，系统在150篇文档上取得了\textbf{92.2\%}的宏平均F1值（Precision 90.5\%，Recall 94.0\%），其中发生时间、经济损失和响应级别三个信息点达到了100\%的准确率。多模态融合对含图片文档的灾害类型识别置信度提升了约21\%，验证了跨模态信息互补的有效性。

未来可在以下方向深入探索：（1）引入预训练语言模型（如BERT-NER）提升文本抽取的精度，特别是对组织名称和复杂地名的识别；（2）采用深度学习图像分类模型（如ResNet、CLIP）替代色彩启发式分析，实现更精确的灾害场景识别；（3）接入实时新闻API实现在线抽取和监测；（4）结合GIS地理信息系统实现灾害事件的空间可视化。


\begin{thebibliography}{12}
\bibitem{jurafsky2024}
    Jurafsky D, Martin J H.
    \textit{Speech and Language Processing}. 3rd ed. 2024.

\bibitem{manning2008}
    Manning C D, Raghavan P, Schütze H.
    \textit{Introduction to Information Retrieval}.
    Cambridge University Press, 2008.

\bibitem{grishman1997}
    Grishman R.
    Information Extraction: Techniques and Challenges.
    \textit{Information Extraction}, Springer, 1997: 10--27.

\bibitem{ahn2006}
    Ahn D.
    The Stages of Event Extraction.
    \textit{Proceedings of the Workshop on Annotating and Reasoning about Time and Events}, ACL, 2006.

\bibitem{chen2015}
    Chen Y, Xu L, Liu K, et al.
    Event Extraction via Dynamic Multi-Pooling Convolutional Neural Networks.
    \textit{Proceedings of ACL}, 2015: 167--176.

\bibitem{li2022}
    Li M, Zareian A, Zeng Q, et al.
    Cross-media Structured Common Space for Multimedia Event Extraction.
    \textit{Proceedings of ACL}, 2020: 2557--2568.

\bibitem{tong2022}
    Tong M, Xu B, Wang S, et al.
    Improving Event Detection via Open-domain Trigger Knowledge.
    \textit{Proceedings of ACL}, 2020: 5887--5897.

\bibitem{jieba}
    jieba 中文分词.
    \url{https://github.com/fxsjy/jieba}.

\bibitem{easyocr}
    EasyOCR: Ready-to-use OCR.
    \url{https://github.com/JaidedAI/EasyOCR}.

\bibitem{flask}
    Flask Web Framework.
    \url{https://flask.palletsprojects.com/}.

\bibitem{visjs}
    vis.js -- Dynamic, browser-based visualization library.
    \url{https://visjs.org/}.

\bibitem{pillow}
    Pillow: Python Imaging Library.
    \url{https://pillow.readthedocs.io/}.
\end{thebibliography}

\end{document}
